% =====================================================================
%  Chương 2 — Hạng mục 2: hai tập dữ liệu chuỗi thực tế
%  Mã nguồn: a06/data_stock.py, a06/data_churn.py · Notebook: 01_du_lieu_amzn, 02_du_lieu_kkbox
% =====================================================================
\chapter{Hai tập dữ liệu chuỗi thực tế}
\label{ch:dulieu}

Bài dùng hai tập dữ liệu thật, đủ lớn và thuộc hai loại bài toán khác nhau. AMZN là chuỗi giá cổ phiếu, một bài hồi
quy: dự báo giá đóng cửa của phiên kế tiếp. KKBox Churn là log nghe nhạc theo ngày của người dùng một dịch vụ trả phí,
một bài phân loại: người dùng này có bỏ dịch vụ hay không. Cả hai đều đưa về tensor $B \times T \times 8$ như
Chương~\ref{ch:lythuyet}, nên cùng một lớp mô hình chạy được trên cả hai.

\section{Giá cổ phiếu AMZN}

\subsection{Nguồn và quy mô}

Dữ liệu ngày của Amazon.com, Inc. (NASDAQ: AMZN) được tải từ Yahoo Finance~\cite{yahooamzn} bằng thư viện
\texttt{yfinance}, gồm \SL{amzn/n_raw} phiên từ \SL{amzn/start} (phiên đầu sau IPO) đến \SL{amzn/end}. Mỗi phiên có
giá mở cửa, cao nhất, thấp nhất, đóng cửa và khối lượng. Giá đã điều chỉnh theo các lần chia tách cổ phiếu, gần nhất là
lần chia 20:1 năm 2022. Tệp CSV được lưu trong kho mã nguồn để lần chạy sau dùng đúng dữ liệu này, không phụ thuộc
việc Yahoo cập nhật lại.

\maNguon{data_stock__add_features}{Thêm đặc trưng tài chính từ giá và khối lượng}

Tám đặc trưng đưa vào mô hình là Close, Open, High, Low, Volume, Return (log-return
$r_t = \ln(P_t/P_{t-1})$), MA20 (trung bình 20 phiên) và RSI14 (chỉ số sức mạnh tương đối của Wilder~\cite{wilder1978}).
MA50 và Volatility20 chỉ dùng để phân tích. Bỏ \SL{amzn/n_raw} $-$ \SL{amzn/n_feat} phiên đầu chưa đủ 50 phiên để tính
MA50, còn lại \SL{amzn/n_feat} phiên. Bảng~\ref{tab:amzn-desc} cho thống kê mô tả.

\begin{table}[H]
  \centering\small
  \caption{Thống kê mô tả của AMZN, \SL{amzn/n_feat} phiên}
  \label{tab:amzn-desc}
  \input{generated/tab_amzn_desc}
\end{table}

\begin{figure}[H]
  \centering
  \hinh{f2_price}
  \caption{Giá đóng cửa AMZN (thang log) với MA20, MA50 và khối lượng; vạch đứt là ranh giới train, val, test}
  \label{fig:price}
\end{figure}

Hình~\ref{fig:price} vẽ giá trên thang log. Trong 27 năm giá tăng hơn một nghìn lần, với hai lần sụt mạnh: bong bóng
dot-com 2000--2001 và đợt điều chỉnh 2022. Thang tuyến tính sẽ ép toàn bộ giai đoạn trước 2010 thành một đường nằm sát
trục hoành.

\subsection{Return, phân phối và biến động}

Hình~\ref{fig:return} xem xét log-return. Return trung bình \SL{amzn/ret/mean} mỗi phiên, độ lệch chuẩn
\SL{amzn/ret/std}. Phiên giảm mạnh nhất là \SL{amzn/ret/min_date} ($r = \SL{amzn/ret/min}$), tăng mạnh nhất là
\SL{amzn/ret/max_date} ($r = \SL{amzn/ret/max}$), đều trong giai đoạn dot-com. Phân phối có đuôi dày: độ nhọn dư
(excess kurtosis) là \SL{amzn/ret/kurt}, trong khi phân phối chuẩn có độ nhọn dư bằng 0. Tỉ lệ phiên lệch quá 4 độ lệch
chuẩn là \SL{amzn/ret/tail4}\,\%, nếu phân phối chuẩn thì chỉ \SL{amzn/ret/tail4_gauss}\,\%. Biến động
(ô c) tụ thành cụm: những giai đoạn dao động mạnh kéo dài nhiều tháng. Ô d cho thấy năm cột giá gần như trùng nhau
(tương quan gần 1). Chúng mang thông tin về mức giá, còn Return và RSI14 mang thông tin về thay đổi.

\begin{figure}[H]
  \centering
  \hinh{f2_return}
  \caption{Log-return, phân phối return so với phân phối chuẩn, biến động 20 phiên và tương quan giữa 8 đặc trưng}
  \label{fig:return}
\end{figure}

\subsection{Tính dừng: kiểm định ADF và tự tương quan}

Kiểm định Augmented Dickey--Fuller~\cite{dickey1979} có giả thuyết không là chuỗi có nghiệm đơn vị (không dừng).
Với giá đóng cửa, thống kê ADF là \SL{amzn/adf/close/stat}, $p = \SL{amzn/adf/close/p}$, không bác bỏ được giả thuyết:
giá không dừng. Với return, thống kê là \SL{amzn/adf/return/statSgn}, $p = \SL{amzn/adf/return/p}$: return dừng.
Hình~\ref{fig:acf} cho cùng kết luận. ACF của giá giảm rất chậm, ở trễ 40 vẫn còn \SL{amzn/acf_c40}. ACF của return
gần như nằm trong dải tin cậy $\pm\SL{amzn/acf_band}$; chỉ \SL{amzn/acf_r_out} trong 40 trễ vượt ra ngoài và đều rất
nhỏ.

\begin{figure}[H]
  \centering
  \hinh{f2_acf}
  \caption{ACF và PACF 40 trễ của giá và của return; dải xám là khoảng tin cậy 95\,\%}
  \label{fig:acf}
\end{figure}

Hai kết quả này quyết định cách đặt bài toán. Giá không dừng, nên mô hình học trên giá 1997--2016 sẽ gặp mức giá hoàn
toàn mới ở tập test: giá lớn nhất của test gấp \SL{amzn/p_max_ratio} lần giá lớn nhất của train. Một mô hình dự báo
thẳng mức giá đã chuẩn hoá Min-Max theo train sẽ phải ngoại suy ra ngoài khoảng $[0, 1]$ mà nó từng thấy. Return thì
dừng, nhưng gần như không có tự tương quan, tức là rất khó dự báo từ quá khứ của chính nó. Đó là điều giả thuyết thị
trường hiệu quả~\cite{fama1970} dự đoán, và Chương~\ref{ch:keras} sẽ thấy các mô hình va vào đúng giới hạn này.

\subsection{Cửa sổ trượt và chống rò rỉ}

Bài dự báo log-return phiên kế tiếp $r_{t+1}$ từ cửa sổ \SL{amzn/window} phiên, rồi quy đổi ra giá
$\hat P_{t+1} = P_t\, e^{\hat r_{t+1}}$ để báo sai số bằng USD. Trong mỗi cửa sổ, năm cột giá chia cho Close của ngày
cuối cửa sổ, nên mô hình chỉ thấy giá tương đối (ví dụ ``MA20 cao hơn giá hiện tại 3\,\%''), không thấy mức giá tuyệt đối.
Volume lấy log.

\maNguon{data_stock__make_windows}{Cửa sổ trượt với nhãn là return phiên kế tiếp}

Tổng cộng \SL{amzn/n_windows} cửa sổ, chia theo thời gian 70/15/15 và không xáo trộn (Hình~\ref{fig:window}). Tập
test là bốn năm cuối, từ \SL{amzn/split/test/start} đến \SL{amzn/split/test/end}. Mọi thống kê chuẩn hoá (trung bình, độ
lệch chuẩn của từng đặc trưng và của nhãn) chỉ tính trên phần train. Notebook 01 kiểm tra bằng \texttt{assert} rằng ngày
cuối của train đứng trước ngày đầu của val, và ngày cuối của val đứng trước ngày đầu của test.

\begin{figure}[H]
  \centering
  \hinh{f2_window}
  \caption{Cửa sổ trượt và cách chia theo thời gian của AMZN}
  \label{fig:window}
\end{figure}

\maNguon{data_stock__WindowScaler}{Chuẩn hoá đặc trưng và nhãn, fit chỉ trên train}

Độ lệch chuẩn của return ở ba phần cũng khác nhau: \SL{amzn/ret_std/train} ở train (có cả giai đoạn dot-com),
\SL{amzn/ret_std/val} ở val và \SL{amzn/ret_std/test} ở test. Mô hình học trên một thị trường biến động mạnh hơn nhiều so
với thị trường nó được đánh giá.

\section{KKBox Churn}

\subsection{Nguồn và định nghĩa nhãn}

KKBox là dịch vụ nghe nhạc trực tuyến lớn ở Đài Loan. Năm 2017, KKBox và hội nghị WSDM 2018 mở một cuộc thi trên
Kaggle~\cite{kkbox2017} với dữ liệu thật của người dùng đã ẩn danh. Bài dùng hai tệp:
\begin{itemize}
  \item \texttt{train\_v2.csv}: \SL{kkbox/n_labeled} người dùng, nhãn \texttt{is\_churn} bằng 1 nếu người dùng
        \emph{không} gia hạn gói thuê bao trong 30 ngày sau khi gói hết hạn vào tháng 3/2017;
  \item \texttt{user\_logs\_v2.csv} (1,4 GB, \SL{kkbox/log_rows} dòng): mỗi dòng là một ngày nghe của một người dùng
        trong tháng 3/2017, gồm số bài nghe dưới 25\,\%, 25--50\,\%, 50--75\,\%, 75--98,5\,\% và trên 98,5\,\% độ dài
        bài (\texttt{num\_25} đến \texttt{num\_100}), số bài khác nhau (\texttt{num\_unq}) và tổng số giây nghe
        (\texttt{total\_secs}).
\end{itemize}

\subsection{Lấy mẫu và dựng chuỗi}

Chỉ \SL{kkbox/n_with_log} người có cả nhãn lẫn ít nhất một ngày log. Từ đó bài lấy mẫu phân tầng
\SL{kkbox/n_sample} người (seed 42). Tệp log không được đọc trọn vào RAM: \texttt{pyarrow} đọc theo khối 64 MB, lượt
thứ nhất lấy tập người có log, lượt thứ hai giữ lại \SL{kkbox/log_rows_kept} dòng của người trong mẫu. Cả quá trình mất
khoảng một phút.

\maNguon{data_churn__sample_users}{Lấy mẫu phân tầng những người vừa có nhãn vừa có log}

Mỗi người dùng thành một ma trận $31 \times 8$: bảy cột log lấy $\log(1+x)$ vì phân phối đuôi rất dài, cột thứ tám là
cờ \texttt{active} (ngày đó có nghe hay không). Ngày không có log để 0. Một ngày có thể có nhiều dòng (nhiều thiết bị)
nên được cộng dồn. \SL{kkbox/secs_clipped} dòng có \texttt{total\_secs} âm hoặc vượt 86.400 giây (dài hơn một ngày) là
lỗi ghi log và được kẹp về $[0, 86.400]$.

\maNguon{data_churn__build_sequences}{Gom log theo ngày thành tensor $(N, 31, 8)$}

Chia 70/15/15 phân tầng theo nhãn: \SL{kkbox/n_train} người train, \SL{kkbox/n_val} val và \SL{kkbox/n_test} test.
Chuẩn hoá z-score từng đặc trưng, thống kê tính trên mọi ngày của mọi người trong train.

\subsection{Phân tích dữ liệu}

Nhãn lệch mạnh: \SL{kkbox/labels/churn} người churn trên \SL{kkbox/labels/total}, tức \SL{kkbox/churn_rate_sample}\,\%
(Hình~\ref{fig:kklabels}). Tỉ lệ này gần như giữ nguyên qua hai bước lọc (\SL{kkbox/churn_rate_labeled}\,\% trên toàn bộ
người có nhãn), nên mẫu đại diện được tập gốc. Một mô hình luôn đoán ``ở lại'' đã đạt accuracy
$100 - \SL{kkbox/churn_rate_sample} \,\%$, vì vậy accuracy không phải thước đo phù hợp; bài dùng F1, Recall, ROC-AUC và
PR-AUC.

\begin{figure}[H]
  \centering
  \hinh{f2_kkbox_labels}
  \caption{Phân bố nhãn trong mẫu và tỉ lệ churn qua các bước lọc}
  \label{fig:kklabels}
\end{figure}

Hình~\ref{fig:kktraj} là kết quả thú vị nhất của phần phân tích. Đầu tháng 3, nhóm sẽ churn nghe \emph{nhiều hơn} nhóm ở
lại ở cả ba đặc trưng. Trong tuần thứ hai (giữa ngày 6 và ngày 11) hai đường cắt nhau, rồi nhóm churn giảm đều đến cuối
tháng, còn nhóm ở lại gần như đi ngang. Tính trung bình cả
tháng, người churn có \SL{kkbox/active_days/churn} ngày nghe, người ở lại có \SL{kkbox/active_days/loyal} ngày. Tín hiệu
nằm ở \emph{xu hướng} theo thời gian chứ không nằm ở mức trung bình, và đó là loại tín hiệu mô hình hồi quy có thể đọc
được mà mô hình trên đặc trưng tổng hợp sẽ bỏ lỡ.

\begin{figure}[H]
  \centering
  \hinh{f2_kkbox_traj}
  \caption{Quỹ đạo trung bình 31 ngày của nhóm ở lại và nhóm churn}
  \label{fig:kktraj}
\end{figure}

\begin{figure}[H]
  \centering
  \begin{minipage}{0.47\textwidth}\centering\hinh{f2_kkbox_hist}\end{minipage}\hfill
  \begin{minipage}{0.5\textwidth}\centering\hinh{f2_kkbox_corr}\end{minipage}
  \caption{Trái: phân phối tổng giờ nghe trong tháng. Phải: tương quan giữa 8 đặc trưng}
  \label{fig:kkhist}
\end{figure}

Hình~\ref{fig:kkhist} (trái) cho thấy phân phối tổng giờ nghe của hai nhóm chồng lên nhau gần hết; trung vị là
\SL{kkbox/hours_med/churn} giờ với nhóm churn và \SL{kkbox/hours_med/loyal} giờ với nhóm ở lại. Các đặc trưng tương quan
mạnh với nhau (phải): \texttt{num\_100} và \texttt{total\_secs} có hệ số \SL{kkbox/corr_100_secs}. Bảng~\ref{tab:ttest}
so hai nhóm bằng kiểm định $t$ của Welch~\cite{welch1947}. Mọi khác biệt đều có ý nghĩa thống kê vì mẫu lớn, nhưng
độ chênh của trung bình nhỏ so với độ phân tán trong từng nhóm. Hành vi nghe chỉ phân biệt được một phần người churn;
phần còn lại phụ thuộc thông tin không có trong log, như phương thức thanh toán hay gói cước. Chương~\ref{ch:pytorch}
sẽ đo mức trần này.

\begin{table}[H]
  \centering\small
  \caption{Kiểm định $t$ Welch giữa nhóm ở lại và nhóm churn, trung bình 31 ngày của từng người}
  \label{tab:ttest}
  \input{generated/tab_ttest}
\end{table}

\maNguon{data_churn__stratified_split}{Chia train/val/test phân tầng theo nhãn}
